\documentclass[10pt,a4paper]{amsart}
\usepackage[margin=19mm]{geometry}
\usepackage{amsmath,amssymb,mathtools}
\usepackage[scr=rsfs,cal=euler]{mathalfa}
\usepackage{xfrac}
\usepackage[unicode,hidelinks,bookmarks=false]{hyperref}
\newcommand{\ie}{i.e.\ }
\newcommand{\cf}{cf.\ }
\newcommand{\resp}{respectively\ }
\newcommand{\Subsection}{\subsection}
\providecommand{\nobreakdash}{\nobreak\hbox{-}}
\setlength{\emergencystretch}{2em}
\multlinegap0pt
\theoremstyle{plain}
\newtheorem{thm}{Theorem}
\newtheorem{dfn}[thm]{Definition}
\newtheorem{lmm}[thm]{Lemma}
\newtheorem{clm}[thm]{Claim}
\newtheorem{cor}[thm]{Corollary}
\newtheorem{cnj}[thm]{Conjecture}
\title[Third eigenvalue of graphs]{Strengthened upper bound on the third eigenvalue of graphs: the block construction $C_6^{[t]}$ is optimal}
\author{Sida Li}
\date{Source: arXiv:2501.07494v2 (CC BY 4.0)}
\begin{document}
\maketitle
\noindent\emph{Source and licence.} Strengthened upper bound on the third eigenvalue of graphs. Version arXiv:2501.07494v2 (CC BY 4.0), distributed under the Creative Commons Attribution 4.0 International licence, \url{http://creativecommons.org/licenses/by/4.0/}, as recorded in the arXiv licence metadata for this exact version and shown on the abstract page \url{https://arxiv.org/abs/2501.07494v2}. Full licence notice: licenses/arxiv-2501.07494-CC-BY-4.0.txt.\par\smallskip

\noindent\textit{Editorial scope.} Counted unit: the article's Theorem with source label \texttt{thm:block} (source0.tex lines 801--803), the sharp bound on the sum of the second and third eigenvalues of the closed vertex multiplication of the six-cycle by the block pattern, stated with its equality case; the article places this theorem in its appendix, where it proves that the block construction used in the proof of the article's main theorem is optimal among the closed vertex multiplications it considers. It is delivered together with the article's complete original proof of it (source0.tex lines 804--830, one \texttt{proof} environment of 27 lines, adjacent to the statement). No mathematics is written, repaired or re-derived: statement and proof are the article's own text line for line, in the article's own notation (the closed vertex multiplication, the divisor matrix of the equitable partition, the characteristic polynomial and its cubic factor, the normalisation of the block parameters, the eigenvalue ordering), and no retained line is substituted, re-flowed or omitted. No further part of the article is required as context and none is retained: the counted proof uses only the objects introduced in its own statement, so no cross-reference and no citation occurs in any retained line; consequently the wrapper carries no bibliography and no bibliography entry, and no context passage is delivered. Nothing else of the article is retained: its main theorem, its claims, its figures and its remaining appendices are omitted. Editorial formatting only: an AMS \texttt{amsart} preamble that restates the article's theorem-like environment declarations, in which Theorem, Definition, Lemma, Claim, Corollary and Conjecture share one counter as in the article, and that adds no mathematical macro, because every control sequence used by the retained lines is standard. The article is typeset in its own \texttt{article}-class format with its own fonts; no class or style file is shipped with this package and the wrapper is an amsart document. Numbering differs from the article: the article resets the shared counter at the start of its appendix and prints the counted theorem as its first appendix theorem, whereas the wrapper keeps a single continuous counter and retains one environment, so the counted theorem prints as Theorem 1. No picture environment and no graphical inclusion occurs in any retained line, and the package delivers no image asset. One bundled source file, \texttt{sources/171555/source0.tex}, accompanies the delivered item as the exact source of every retained span. Every retained span is recorded in \texttt{delivery-evidence.json} with method \texttt{exact\_tex}.
\begin{thm}\label{thm:block}
    Suppose $G$ is the closed vertex multiplication of $C_6$ by $[a, b, c, a, b, c]$. Then $\lambda_2 + \lambda_3 \le \frac{2n}{3} - 2$, with equality iff $a = b = c$. 
\end{thm}

\begin{proof}
    Add a self-loop at each vertex and take the divisor matrix for the equitable partition. 
    \[\begin{pmatrix}
        a & b & 0 & 0 & 0 & c \\
        a & b & c & 0 & 0 & 0 \\
        0 & b & c & a & 0 & 0 \\
        0 & 0 & c & a & b & 0 \\
        0 & 0 & 0 & a & b & c \\
        a & 0 & 0 & 0 & b & c
    \end{pmatrix}.\]
    
    This has characteristic polynomial $x^2(x-(a+b+c))P(x)$ where
    \[P(x) = x^3 - (a+b+c)x^2 + 4abc.\]
    
    Thus, the eigenvalues of $G$ are 0, $a+b+c$ which corresponds to the regularity, and the roots of $P(x) = 0$. Normalising $a + b + c = \frac12$, we get $x^3 - \frac12 x^2 + 4\alpha$ where $0 \le \alpha \le \frac1{216}$ due to AM-GM. Note that $P(1/3) = -\frac1{54} + 4\alpha \le 0$, hence it either lies between the first two roots or before the last root. Yet $P(0) = 4\alpha \ge 0$ hence the last root is non-positive. 

    If $\frac13 + t$ is a root for $P$ with $0 \le t \le \frac13$, then
    \begin{align*}
        0 &= \left(\frac{1}{27} + \frac13 t + t^2 + t^3\right) - \frac12 \left(\frac19 + \frac23 t + t^2\right) - 4\alpha \\ 
        \Rightarrow -\frac23 t - 2t^3 + \frac23 t &= \left(\frac{1}{27} - \frac13 t + t^2 - t^3\right) - \frac12 \left(\frac19 - \frac23 t + t^2\right) - 4\alpha \\
        \Rightarrow P\left(\frac13 - t\right) &= -2t^3 \le 0.
    \end{align*}
    
    Thus, $\frac13 - t$ lies between the first two roots. In particular, $\frac13 + t$ must be the first root, and the sum of the first two roots is at most $\frac23$, with equality iff $t = 0$ iff $\frac13$ is a root. This happens at $abc = \frac1{216}$ and hence $a = b = c = \frac16$ as the equality case for AM-GM. 

    Since we obtained the eigenvalues of $A(G) + I$ from $P$, we have $(\lambda_2 + 1) + (\lambda_3 + 1) \le \frac{2n}{3}$ with equality iff $a = b = c$. 
\end{proof}

\end{document}
